% AUDIT DE VERACITE 2026-09-20 -- corrections des citations de rapport de jury.
% Deux occurrences de la meme citation fabriquee (« les fonctions explicites doivent etre nommees et
% justifiees », attribuee au jury 2025) : verbatim introuvable dans les 8 rapports 2018-2025, et la
% locution « fonctions explicites » n'y figure nulle part.
%   - Bloc 3, corrige de QT2 : la note « Le jury note : ... » est SUPPRIMEE ; le conseil de methode
%     (justifier par le sens du flux d'energie) est conserve, sans attribution au jury.
%   - Bloc 4, encadre jury : la citation 2025 est REMPLACEE par le verbatim authentique du
%     Rapport_Jury_2025.md ligne 486 (texte identique au Rapport_Jury_2024.md ligne 515), qui porte
%     bien sur l'exigence de justification des resultats.
% La citation « Rapport jury 2024 (B13) » sur la MLI et le THD n'a pas ete modifiee.
\expandafter\ifx\csname CAPETINCLUDEMODE\endcsname\relax
\documentclass[11pt,a4paper]{article}
% preamble.tex — packages et commandes communes a toutes les fiches CAPET SII IE
\input{glyphtounicode}
\pdfgentounicode=1
\usepackage{cmap}
\usepackage[T1]{fontenc}
\usepackage[utf8]{inputenc}
\usepackage[french]{babel}
% babel[french] active a \begin{document} des raccourcis de ponctuation (:;!?)
% qui cassent la syntaxe to[...] de circuitikz (erreur "+ or - expected") ;
% on les desactive, ce qui ne change que l'espacement fin avant ces signes.
\AtBeginDocument{\shorthandoff{:;!?}}
\usepackage[a4paper, top=1.6cm, bottom=1.6cm, left=1.4cm, right=1.4cm]{geometry}
\setlength{\headheight}{14.5pt}
\usepackage{amsmath, amssymb, amsfonts}
\usepackage{siunitx}
\sisetup{locale=FR, output-decimal-marker={,}, inter-unit-product=\ensuremath{{}\cdot{}},
         group-minimum-digits=4}
% Unites non fournies par siunitx, declarees ici une fois pour toutes (fiches + exercices).
% Les gardes \ifdefined evitent tout conflit avec une redeclaration locale restee dans un
% fichier : ne pas retirer les gardes, ne pas redeclarer ces unites dans les fichiers.
\DeclareSIUnit{\jour}{j}
\ifdefined\var\else\DeclareSIUnit{\var}{var}\fi   % puissance reactive
\ifdefined\vad\else\DeclareSIUnit{\vad}{VAD}\fi   % puissance deformante
\ifdefined\tr\else\DeclareSIUnit{\tr}{tr}\fi      % tour (tr/min)
\usepackage{array, booktabs, multirow, tabularx}
\usepackage{xcolor}
\usepackage[most, breakable]{tcolorbox}
\usepackage{multicol}
\usepackage{enumitem}
\usepackage{mathtools}
\usepackage{lmodern}
% Refonte visuelle 2026 (direction "moderne coloré") : Lato pour le texte,
% maths sans-serif assorties via newtxsf (charge apres amsmath/amssymb).
\usepackage[default]{lato}
\usepackage{newtxsf}
\usepackage{microtype}
\usepackage{fancyhdr}
\usepackage{tikz}
\usetikzlibrary{positioning,arrows.meta}
\usepackage{circuitikz}
\usepackage{lastpage}

% ── Couleurs par domaine (palette modernisée 2026) ────────────────────
\definecolor{colElec}{RGB}{29, 78, 216}        % Électrotechnique — bleu vif
\definecolor{colElecLight}{RGB}{219, 234, 254} % fond bleu clair
\definecolor{colAuto}{RGB}{21, 128, 61}        % Automatique — vert
\definecolor{colMeca}{RGB}{194, 65, 12}        % Mécanique — orange
\definecolor{colTherm}{RGB}{185, 28, 28}       % Thermique — rouge
\definecolor{colQT}{RGB}{232, 244, 233}        % (hérité, plus utilisé par questiontype)
\definecolor{colJury}{RGB}{255, 248, 225}      % fond jury (crème)
\definecolor{colJuryFrame}{RGB}{202, 138, 4}   % bordure/bandeau jury (ambre)
\definecolor{colGrayLight}{RGB}{245, 245, 247} % fond gris neutre
\definecolor{colSeparator}{RGB}{180, 180, 180} % gris séparateur
\definecolor{colExo}{RGB}{8, 145, 178}          % Banque d'exercices — cyan (distinct des fiches)
% ── Encre et gris de la refonte ───────────────────────────────────────
\definecolor{colInk}{RGB}{23, 32, 42}          % texte courant (quasi-noir)
\definecolor{colSlate}{RGB}{83, 98, 118}       % gris ardoise (méta, en-têtes)
\definecolor{colRule}{RGB}{203, 213, 225}      % filets clairs
\definecolor{colJuryBg}{RGB}{255, 248, 225}    % fond crème jury
\definecolor{colPresence}{RGB}{83, 98, 118}    % bloc 5 (présence dans les sujets)
\AtBeginDocument{\color{colInk}}

% ── Style encadré exercice + corrigé (banque d'exercices) ─────────────
\tcbset{
  exoBox/.style={
    breakable,
    enhanced,
    colback=colExo!8!white,
    colframe=colExo,
    colbacktitle=colExo,
    coltitle=white,
    fonttitle=\bfseries\small,
    left=4pt, right=4pt, top=3pt, bottom=3pt,
    before skip=5pt, after skip=5pt,
  }
}

% ── Bascule enonce / corrige (banque d'exercices) ─────────────────────
% \CAPETCORRIGE est defini en ligne de commande (pdflatex "\def\CAPETCORRIGE{1}\input{...}")
% pour compiler la variante corriges d'un meme fichier source ; par defaut (non defini),
% la compilation produit la variante enonces. \CAPETSHOW{enonce}{corrige} choisit la branche.
%
% TROIS MODES (refonte 2026-09-20, manuel de revision) :
%   - aucun drapeau   -> imprime l'enonce            (variante enonces, INCHANGE)
%   - \CAPETCORRIGE   -> imprime le corrige          (variante corriges, INCHANGE)
%   - \CAPETMANUEL    -> imprime l'enonce ET met le corrige de cote pour l'annexe
%
% Le mode manuel stocke chaque corrige dans une macro (\csgdef), sans ecriture
% fichier ni re-tokenisation : le contenu (circuitikz, accents, %) traverse intact.
% L'annexe les rejoue ensuite dans l'ordre via \CAPETANNEXE.
\usepackage{etoolbox}

\newcounter{CAPETexo}          % numero continu d'exercice dans le manuel
\setcounter{CAPETexo}{0}
\newcounter{CAPETannexe}       % compteur de relecture en annexe

% Titre courant de la notion en cours (pose par \CAPETNOTION), sert d'intertitre annexe
\newcommand{\CAPETnotioncode}{}
\newcommand{\CAPETnotiontitre}{}

% \CAPETNOTION{code}{titre} : ouvre une notion dans le manuel (memorise le libelle,
% et marque le rang du premier exercice de la notion pour l'intertitre d'annexe).
\newcommand{\CAPETNOTION}[2]{%
  \renewcommand{\CAPETnotioncode}{#1}%
  \renewcommand{\CAPETnotiontitre}{#2}%
  \csgdef{CAPETnotionstart\the\numexpr\value{CAPETexo}+1\relax}{#1 -- #2}%
}

% ── \degree robuste ───────────────────────────────────────────────────
% Certains fichiers d'exercices utilisent \degree. Ni le preambule ni siunitx
% ne le definissent, et \providecommand local ne suffit pas : sous hyperref,
% \degree existe deja comme commande factice (\ERROR), donc \providecommand
% s'abstient et l'appel echoue a la compilation du manuel (erreur invisible
% en compilation isolee). On le (re)definit donc inconditionnellement ici,
% apres le chargement des packages, pour tout le corpus.
\AtBeginDocument{\providecommand{\degree}{}\renewcommand{\degree}{\ensuremath{^\circ}}}

% Libelle decoratif : les fichiers d'exercices utilisent aussi \CAPETSHOW dans
% leur \fancyhead et leur bloc de titre, pour afficher "(enonces)" ou "(corriges)".
% Ces appels-la ne sont PAS des exercices : en mode manuel ils ne doivent ni
% consommer un numero, ni alimenter l'annexe. On les reconnait a leur premier
% argument, strictement egal au libelle ci-dessous.
\newcommand{\CAPETLIBELLEENONCES}{(\'enonc\'es)}

\ifdefined\CAPETMANUEL
  % --- Mode manuel : enonce en place, corrige differe vers l'annexe ---
  \newcommand{\CAPETSHOW}[2]{%
    \def\CAPETargun{#1}%
    \ifx\CAPETargun\CAPETLIBELLEENONCES
      #1%            appel decoratif : rendu tel quel, hors numerotation
    \else
      \CAPETSHOWEXO{#1}{#2}%
    \fi
  }
  \newcommand{\CAPETSHOWEXO}[2]{%
    \stepcounter{CAPETexo}%
    \csgdef{CAPETcorrige\theCAPETexo}{#2}%
    \edef\CAPETtmpnum{\theCAPETexo}%
    #1%
    \par\vspace{2pt}%
    {\footnotesize\itshape\color{colExo}%
     \hyperref[capet:cor:\CAPETtmpnum]{Corrige de l'exercice \CAPETtmpnum\ $\rightarrow$ p.~\pageref{capet:cor:\CAPETtmpnum}}%
     \label{capet:eno:\CAPETtmpnum}}%
  }
  % \CAPETANNEXE : rejoue tous les corriges captures, dans l'ordre, avec intertitres
  \newcommand{\CAPETANNEXE}{%
    \setcounter{CAPETannexe}{0}%
    \loop\ifnum\value{CAPETannexe}<\value{CAPETexo}%
      \stepcounter{CAPETannexe}%
      \edef\CAPETcurnum{\theCAPETannexe}%
      \ifcsname CAPETnotionstart\CAPETcurnum\endcsname
        \par\medskip
        {\normalsize\bfseries\color{colExo}\csuse{CAPETnotionstart\CAPETcurnum}}%
        \par\nobreak\smallskip
      \fi
      \begin{tcolorbox}[exoBox, breakable,
        title={Corrige de l'exercice \CAPETcurnum\
               \hfill\normalfont\footnotesize
               \hyperref[capet:eno:\CAPETcurnum]{enonce p.~\pageref{capet:eno:\CAPETcurnum}}}]%
      \label{capet:cor:\CAPETcurnum}%
      \csuse{CAPETcorrige\CAPETcurnum}%
      \end{tcolorbox}%
    \repeat
  }
\else
  % --- Modes historiques, comportement strictement inchange ---
  \newcommand{\CAPETSHOW}[2]{\ifdefined\CAPETCORRIGE #2\else #1\fi}
  \newcommand{\CAPETANNEXE}{}
\fi

% ── Difficulté 1 à 3 (étoiles pleines, amssymb) ───────────────────────
\newcommand{\difficulte}[1]{%
  \ifcase#1\relax
  \or $\bigstar$\;niveau 1/3%
  \or $\bigstar\bigstar$\;niveau 2/3%
  \or $\bigstar\bigstar\bigstar$\;niveau 3/3%
  \fi
}

% ── Style encadré extrait de document technique (banque d'exercices) ──
\tcbset{
  docTech/.style={
    breakable,
    enhanced,
    colback=colGrayLight,
    colframe=colSeparator,
    fonttitle=\bfseries\small\itshape,
    coltitle=black,
    colbacktitle=colGrayLight,
    left=4pt, right=4pt, top=3pt, bottom=3pt,
    before skip=4pt, after skip=4pt,
  }
}

% ── Style tcolorbox de section principale ────────────────────────────
% Refonte 2026 : bandeau titre coloré arrondi + fond très pâle, sans cadre.
\tcbset{
  ficheSection/.style={
    breakable, enhanced jigsaw,
    frame hidden, boxrule=0pt, arc=1.8mm,
    colback=#1!4!white,
    colbacktitle=#1,
    coltitle=white,
    fonttitle=\sffamily\bfseries,
    left=7pt, right=7pt, top=5pt, bottom=5pt,
    toptitle=3pt, bottomtitle=3pt,
    before skip=10pt, after skip=8pt,
  }
}

% ── Style encadré question type ──────────────────────────────────────
% Carte blanche, bandeau sombre, liseré très clair.
\tcbset{
  questiontype/.style={
    breakable, enhanced jigsaw,
    colback=white,
    frame hidden, boxrule=0pt, arc=1.8mm,
    colbacktitle=colElec!78!black,
    coltitle=white,
    fonttitle=\sffamily\bfseries\small,
    left=7pt, right=7pt, top=4pt, bottom=4pt,
    toptitle=2.5pt, bottomtitle=2.5pt,
    before skip=8pt, after skip=8pt,
    borderline={0.6pt}{0pt}{colElec!12!white},
  }
}

% ── Style encadré jury ───────────────────────────────────────────────
% Bandeau ambre, fond crème, sans cadre.
\tcbset{
  jury/.style={
    breakable, enhanced jigsaw,
    colback=colJuryBg,
    frame hidden, boxrule=0pt, arc=1.8mm,
    colbacktitle=colJuryFrame,
    coltitle=white,
    fonttitle=\sffamily\bfseries,
    left=7pt, right=7pt, top=4pt, bottom=4pt,
    toptitle=3pt, bottomtitle=3pt,
    before skip=10pt, after skip=8pt,
  }
}

% ── En-tête / pied de page communs (refonte 2026) ─────────────────────
% #1 code fiche, #2 titre court, #3 couleur domaine
\newcommand{\ficheHeaders}[3]{%
  \fancyhead[L]{\small\sffamily{\bfseries\textcolor{#3}{#1}}\hspace{0.7em}\textcolor{colSlate}{#2}}%
  \fancyhead[R]{\small\sffamily\textcolor{colSlate}{CAPET SII IE\, --- \,Session 2026}}%
  \fancyfoot[C]{\small\sffamily\textcolor{colSlate}{\thepage\,/\,\pageref{LastPage}}}%
  \fancyhead[C]{}\fancyfoot[L]{}\fancyfoot[R]{}%
}

% ── Bloc en-tête de fiche (BLOC 0, refonte 2026) ──────────────────────
% #1 code  #2 titre  #3 couleur domaine  #4 domaine  #5 ligne présence  #6 sous-thèmes
\newcommand{\ficheEntete}[6]{%
  \begin{tcolorbox}[enhanced, breakable,
    frame hidden, boxrule=0pt, arc=2.5mm,
    interior style={left color=#3, right color=#3!75!black},
    left=9pt, right=9pt, top=7pt, bottom=7pt,
    before skip=0pt, after skip=6pt]
    {\color{white}\sffamily\bfseries\LARGE #1\quad #2}\\[4pt]
    {\color{#3!22!white}\sffamily\small \MakeUppercase{#4}\quad\textbar\quad #5}\\[1pt]
    {\color{#3!35!white}\sffamily\footnotesize Sous-thèmes : #6}
  \end{tcolorbox}%
}

% ── Mise en page ─────────────────────────────────────────────────────
\pagestyle{fancy}
\fancyhf{}
\renewcommand{\headrulewidth}{0.5pt}
\renewcommand{\footrulewidth}{0pt}
\AtBeginDocument{%
  \renewcommand{\headrule}{{\color{colRule}\hrule width\headwidth height 0.5pt}}%
}

\setlist[itemize]{leftmargin=1.2em, itemsep=1pt, topsep=2pt, parsep=0pt}
\setlist[enumerate]{leftmargin=1.2em, itemsep=3pt, topsep=2pt, parsep=0pt}

\setlength{\parskip}{2pt}
\setlength{\columnsep}{10pt}

\def\CAPETBEGINFILE{\begin{document}}
\def\CAPETENDFILE{\end{document}}
\else
\def\CAPETBEGINFILE{}
\def\CAPETENDFILE{}
\fi

\ficheHeaders{T09}{Onduleurs}{colElec}

\CAPETBEGINFILE

% -----------------------------------------------------------------------
%  BLOC EN-TETE
% -----------------------------------------------------------------------

\ficheEntete{T09}{Onduleurs (monophasé et triphasé)}{colElec}%
  {électrotechnique - électronique de puissance}%
  {\textbf{Présence :} \textbf{12/13} sujets (2013-2026) - thème \emph{très fréquent}}%
  {pleine onde (créneau) - MLI sinusoïdale - THD - norme IEC 61000-2-2 - onduleur triphasé (machine asynchrone, MSAP) - îlotage vs réseau}

% -----------------------------------------------------------------------
%  1. RAPPELS DE COURS
% -----------------------------------------------------------------------

\begin{tcolorbox}[
  ficheSection=colElec,
  title={1. Rappels de cours},
]

% -- 1.1 Onduleur pleine onde -----------------------------------------
\textbf{1.1 Onduleur pleine onde - tension en créneau}

\smallskip
Un créneau $\pm V_{DC}$ a une valeur efficace égale à $V_{DC}$ (le carré est constant) :
\[
  V_{eff,mono} = V_{DC}
\]

\smallskip
\begin{center}
\begin{tikzpicture}[scale=0.85, font=\footnotesize]
  \draw[->] (0,-1.6) -- (0,1.6) node[above] {$v(t)$};
  \draw[->] (-0.3,0) -- (6.3,0) node[right] {$t$};
  \draw[thick, colElec] (0,1.2) -- (3,1.2) -- (3,-1.2) -- (6,-1.2);
  \node[left] at (-0.05,1.2) {\scriptsize $+V_{DC}$};
  \node[left] at (-0.05,-1.2) {\scriptsize $-V_{DC}$};
  \draw[dotted] (3,-1.2) -- (3,1.2);
  \draw[dashed] (6,0) -- (6,-1.8);
  \node[below] at (6,-1.8) {$T$};
\end{tikzpicture}
\end{center}

Décomposition en série de Fourier (harmoniques impairs seulement) :
\[
  V_n = \frac{4V_{DC}}{n\pi} \quad (n \text{ impair} : 1, 3, 5, 7, \ldots)
  \qquad;\qquad
  V_1 = \frac{4V_{DC}}{\pi} \approx 1{,}273\,V_{DC}
\]
Pour un onduleur triphasé équilibré, les rangs multiples de 3 disparaissent des tensions composées : rangs présents = $6k \pm 1$ ($1, 5, 7, 11, 13, \ldots$).

Valeur efficace tension composée (pleine onde triphasé, commande $120^\circ$) :
\[
  V_{eff,comp} = \sqrt{\frac{2}{3}} \cdot V_{DC} \approx 0{,}816\,V_{DC}
\]

\medskip
\hrule\vspace{4pt}

% -- 1.2 Taux de distorsion harmonique --------------------------------
\textbf{1.2 Taux de distorsion harmonique (THD)}

\smallskip
\[
  \mathrm{THD} = \frac{\sqrt{\displaystyle\sum_{n=2}^{\infty} V_n^2}}{V_1}
  \qquad \text{(dénominateur = fondamental } V_1 \text{, pas la valeur efficace totale)}
\]

\begin{center}
\renewcommand{\arraystretch}{1.35}
\begin{tabular}{llll}
\toprule
\textbf{Structure} & \textbf{THD} & \textbf{Norme IEC 61000-2-2} & \textbf{Conforme ?} \\
\midrule
Pleine onde triphasé & $61\,\%$ & $< 25\,\%$ & Non \\
MLI ($f_p = 6\,\mathrm{kHz}$) & $25\,\%$ & $< 25\,\%$ & Limite \\
\bottomrule
\end{tabular}
\end{center}

\medskip
\hrule\vspace{4pt}

% -- 1.3 MLI sinusoedale ----------------------------------------------
\textbf{1.3 MLI sinusoïdale - indices de modulation}

\smallskip
\begin{center}
\renewcommand{\arraystretch}{1.35}
\begin{tabular}{lll}
\toprule
\textbf{Indice} & \textbf{Définition} & \textbf{Rôle} \\
\midrule
Amplitude : $m_a = \hat{V}_{ref}/\hat{V}_{porteuse}$ & $m_a \leq 1$ : linéaire & Contrôle l'amplitude du fondamental \\[4pt]
Fréquence : $m_f = f_{porteuse}/f_{fondamental}$ & Ex. $m_f = 6000/50 = 120$ & Localise les harmoniques à $m_f \cdot f$ \\[4pt]
Fondamental MLI : $\hat{V}_1 = m_a \cdot V_{DC}/2$ & (monophasé) & Valide pour $m_a \leq 1$ \\
\bottomrule
\end{tabular}
\end{center}

\smallskip
La MLI repousse les harmoniques à haute fréquence ($\approx m_f \cdot f = 6\,\mathrm{kHz}$), naturellement atténuées par l'effet inductif des enroulements de la machine.

\medskip
\hrule\vspace{4pt}

% -- 1.4 Frequence statorique (machine) -------------------------------
\textbf{1.4 Fréquence de commande - machine tournante}

\smallskip
\[
  f_s = \frac{p \cdot n_m}{60} \qquad (p = \text{nb paires de pôles},\; n_m = \text{vitesse en tr/min})
\]
L'onduleur doit générer \textbf{exactement} cette fréquence statorique pour imposer la vitesse souhaitée.

\end{tcolorbox}

% -----------------------------------------------------------------------
%  2. FORMULAIRE SYNTHETIQUE
% -----------------------------------------------------------------------

\begin{tcolorbox}[
  ficheSection=colElec,
  title={2. Formulaire synthétique},
]
\small
\begin{multicols}{2}
\renewcommand{\arraystretch}{1.45}
\begin{tabular}{@{}p{4.5cm}ll@{}}
\toprule
\textbf{Formule} & \textbf{Usage} & \textbf{U.} \\
\midrule
\multicolumn{3}{@{}l}{\footnotesize\textit{Pleine onde}} \\
$V_{eff} = V_{DC}$ & Créneau mono & V \\
$V_{eff,comp} = \sqrt{2/3}\,V_{DC}$ & Triphasé comp. & V \\
$V_n = 4V_{DC}/(n\pi)$ & Harmonique $n$ & V \\
$V_1 \approx 1{,}273\,V_{DC}$ & Fondamental & V \\
\midrule
\multicolumn{3}{@{}l}{\footnotesize\textit{THD}} \\
$\mathrm{THD} = \sqrt{\sum V_n^2}/V_1$ & Général & - \\
$\mathrm{THD}_{po,mono} \approx 0{,}484$ & Pleine onde & - \\
Norme IEC : THD $< 25\,\%$ & Limite & - \\
\bottomrule
\end{tabular}

\columnbreak

\renewcommand{\arraystretch}{1.45}
\begin{tabular}{@{}p{4.5cm}ll@{}}
\toprule
\textbf{Formule} & \textbf{Usage} & \textbf{U.} \\
\midrule
\multicolumn{3}{@{}l}{\footnotesize\textit{MLI}} \\
$m_a = \hat{V}_{ref}/\hat{V}_{porteuse}$ & Indice ampli & - \\
$m_f = f_p/f$ & Indice fréq. & - \\
$\hat{V}_1 = m_a \cdot V_{DC}/2$ & Fondamental & V \\
\midrule
\multicolumn{3}{@{}l}{\footnotesize\textit{Machine tournante}} \\
$f_s = p \cdot n_m/60$ & Fréq. statoriq. & Hz \\
$U = \sqrt{3}\,V$ & Comp./simple & V \\
$S = V_{eff} I_{eff}$ & Puiss. app. & VA \\
\midrule
\multicolumn{3}{@{}l}{\footnotesize\textit{Données numériques}} \\
THD pleine onde trip. & $61\,\%$ & - \\
THD MLI $6\,\mathrm{kHz}$ & $25\,\%$ & - \\
\bottomrule
\end{tabular}
\end{multicols}
\normalsize
\end{tcolorbox}

% -----------------------------------------------------------------------
%  3. QUESTIONS TYPES
% -----------------------------------------------------------------------

\begin{tcolorbox}[
  ficheSection=colElec,
  title={3. Questions types - extraites des sujets réels},
]

% -- QT1 -------------------------------------------------------------

\begin{tcolorbox}[questiontype,
  title={QT1 - THD pleine onde vs MLI - conformité norme IEC \hfill\normalfont\footnotesize E2 disciplinaire 2024 (B11-B13)}
]
\textbf{Énoncé.}
Deux cas de commande pour un onduleur triphasé (variateur d'extracteur d'air) :
\begin{itemize}
  \item Pleine onde : $V_1 = 170\,\mathrm{V}$, $V_5 = 95\,\mathrm{V}$, $V_7 = 27\,\mathrm{V}$, $V_{11} = 27\,\mathrm{V}$, $V_{13} = 20\,\mathrm{V}$.
  \item MLI $f_p = 6\,\mathrm{kHz}$ : $V_1 = 170\,\mathrm{V}$, $V_{119} = 37\,\mathrm{V}$, $V_{121} = 22\,\mathrm{V}$.
\end{itemize}
Calculer le THD dans chaque cas. Conclure par rapport à la norme IEC/EN 61000-2-2 ($\mathrm{THD} < 25\,\%$).

\tcbline

\textbf{Corrigé.}

\textbf{Étape 1 — Pleine onde.}
\[
  \mathrm{THD}_1 = \frac{\sqrt{95^2 + 27^2 + 27^2 + 20^2}}{170}
  = \frac{\sqrt{9\,025 + 729 + 729 + 400}}{170}
  = \frac{\sqrt{10\,883}}{170}
  = \frac{104{,}3}{170} \approx \mathbf{0{,}61 = 61\,\%}
\]
Non conforme ($61\,\% > 25\,\%$). \textit{Remarque : le corrigé officiel 2024 affiche $0{,}59$, mais le recalcul exact avec les valeurs du DT3 donne $0{,}61$ — la conclusion de non-conformité est inchangée.}

\textbf{Étape 2 — MLI $6\,\mathrm{kHz}$.}
\[
  \mathrm{THD}_2 = \frac{\sqrt{37^2 + 22^2}}{170}
  = \frac{\sqrt{1\,369 + 484}}{170}
  = \frac{\sqrt{1\,853}}{170}
  \approx \mathbf{0{,}25 = 25\,\%}
\]
Conforme (limite atteinte). La MLI repousse les harmoniques aux rangs $119$ et $121$ (${\approx}\,6\,\mathrm{kHz}$), naturellement filtrés par l'inductance des enroulements de la machine.

\textbf{Conclusion.} Seule la commande MLI atteint la conformité IEC/EN 61000-2-2 : c'est la commande à retenir pour ce variateur, la pleine onde étant à écarter pour un raccordement réseau normal.
\end{tcolorbox}

\vspace{4pt}

% -- QT2 -------------------------------------------------------------

\begin{tcolorbox}[questiontype,
  title={QT2 - Identification des convertisseurs dans une chaîne PV \hfill\normalfont\footnotesize Rapport jury 2025 (Q14)}
]
\textbf{Énoncé.}
Dans la chaîne de conversion d'un site isolé de Mafate : PV - CV1 - bus DC - CV3 - réseau local $230\,\mathrm{V}$, $50\,\mathrm{Hz}$. Identifier la nature de CV1 et CV3, et justifier.

\tcbline

\textbf{Corrigé.}

\textbf{CV1 :} DC - DC, adapte la tension PV variable au bus DC - \textbf{Hacheur Boost} (MPPT).\\
\textbf{CV3 :} DC - AC, convertit le bus DC en tension AC $230\,\mathrm{V}$, $50\,\mathrm{Hz}$ pour le réseau local - \textbf{Onduleur}.

\textit{Justifier par le sens du flux d'énergie (DC→AC = onduleur) plutôt que par le seul schéma.}

\textbf{Conclusion.} Identifier CV1 en hacheur Boost et CV3 en onduleur permet de justifier le choix technologique par le sens de conversion d'énergie exigé par le cahier des charges (PV variable → bus DC stable → réseau local en 230\,V/50\,Hz).
\end{tcolorbox}

\vspace{4pt}

% -- QT3 -------------------------------------------------------------

\begin{tcolorbox}[questiontype,
  title={QT3 - Onduleur triphasé - fréquence statorique (MSAP Urbanloop) \hfill\normalfont\footnotesize E2 appliquée 2024}
]
\textbf{Énoncé.}
Une MSAP ($p = 14$ paires de pôles, rayon roue $R_{roue} = \SI{0.20}{\metre}$, donnée sujet) est alimentée par un onduleur triphasé ($V_{DC} = \SI{72}{\volt}$). La capsule roule à $V_0 = \SI{60}{\kilo\metre\per\hour}$. Calculer la vitesse de rotation $n_m$ puis la fréquence statorique $f_s$ à imposer à l'onduleur.

\tcbline

\textbf{Corrigé.}

\textbf{Étape 1 — Vitesse de rotation.}
\[
  V_0 = \SI{60}{\kilo\metre\per\hour} = \SI{16.67}{\metre\per\second}
  \qquad
  n_m = \frac{V_0}{2\pi R_{roue}} \times 60 = \frac{16{,}67}{2\pi \times 0{,}20} \times 60 \approx 796\,\mathrm{tr/min}
\]

\textbf{Étape 2 — Fréquence statorique.}
\[
  f_s = \frac{p \cdot n_m}{60} = \frac{14 \times 796}{60} \approx \mathbf{\SI{186}{\hertz}}
\]

\textbf{Conclusion.} L'onduleur doit générer exactement $f_s \approx \SI{186}{\hertz}$ pour imposer la vitesse de $\SI{60}{\kilo\metre\per\hour}$ à la MSAP : toute autre fréquence produirait une vitesse différente de la cible du cahier des charges.
\end{tcolorbox}

\end{tcolorbox}

% -----------------------------------------------------------------------
%  4. ATTENTES DU JURY
% -----------------------------------------------------------------------

\begin{tcolorbox}[
  jury,
  title={4. Attentes du jury - Erreurs récurrentes et points de vigilance},
]

\begin{enumerate}[label=\textbf{\arabic*.}, leftmargin=*, itemsep=5pt]

  \item \textbf{$V_{eff}$ pleine onde = $V_{DC}$, pas $V_{DC}/\sqrt{2}$.}
  La tension en créneau $\pm V_{DC}$ a une valeur efficace égale à $V_{DC}$ (le carré vaut toujours $V_{DC}^2$).
  Confondre avec une sinusoïde ($V_{eff} = V_{DC}/\sqrt{2}$) est l'erreur la plus fréquente.

  \item \textbf{THD : dénominateur = fondamental $V_1$, pas la valeur efficace totale.}
  $\mathrm{THD} = \sqrt{\sum V_n^2}/V_1$. Diviser par $V_{eff}$ total donne une valeur fausse.

  \item \textbf{Justifier physiquement le filtrage par la machine.}
  Pas seulement citer que la MLI est meilleure : les harmoniques HF ($6\,\mathrm{kHz}$) sont atténuées par l'effet inductif ($Z = j\omega L$ très élevé à haute fréquence).

  \item \textbf{Harmoniques triphasés : rangs absents des tensions composées.}
  Pour un onduleur triphasé équilibré, les rangs multiples de 3 (3, 9, 15\ldots) disparaissent. Rangs présents : $6k \pm 1$.

  \item \textbf{Fréquence statorique = $p \cdot n_m/60$.}
  Ne pas confondre $p$ (paires de pôles) avec le nombre de pôles. Une machine à $p = 14$ paires a $28$ pôles.

  \item \textbf{Identifier la nature du convertisseur depuis le schéma.}
  DC→AC = onduleur ; AC→DC = redresseur ; DC→DC = hacheur. Nommer et justifier par le sens du flux.

  \item \textbf{Onduleur en îlotage vs connecté réseau.}
  En îlotage, l'onduleur fixe lui-même $f = 50\,\mathrm{Hz}$ et $V = 230\,\mathrm{V}$. En mode réseau, il est synchronisé.

\end{enumerate}

\medskip
\begin{tcolorbox}[
  colback=colJuryFrame!12!colJury,
  colframe=colJuryFrame,
  left=4pt, right=4pt, top=2pt, bottom=2pt,
  fontupper=\small\itshape
]
Rapport jury 2024 (B13) : «~La MLI permet de limiter le THD du courant et donc d'en limiter les ondulations. [\ldots] De plus la suppression des harmoniques en haute fréquence est faisable par l'effet inductif des enroulements de la machine.~»
\textnormal{\footnotesize (Le corrigé officiel B10 affiche THD pleine onde $= 0{,}59$ et THD MLI $= 0{,}25$ ; le recalcul exact du cas pleine onde donne $0{,}61$ — la conclusion de non-conformité de la pleine onde reste identique.)}\\
Rapport jury 2025 : «~Le jury attend des réponses détaillées et argumentées. Des résultats donnés directement, sans calcul et sans justification ne répondent pas aux attendus de l'épreuve.~»
\end{tcolorbox}

\end{tcolorbox}

% -----------------------------------------------------------------------
%  PIED - PRESENCE DANS LES SUJETS
% -----------------------------------------------------------------------

\begin{tcolorbox}[
  ficheSection=colPresence,
  title={Présence dans les sujets analysés (2013-2026)},
]
\begin{center}
\footnotesize\renewcommand{\arraystretch}{1.15}
\begin{tabular}{lccccccccccccc}
\toprule
\textbf{Année} & 13 & 14 & 15 & 16 & 17 & 18 & 19 & 20 & 21 & 22 & 23 & 24 & 25 \\
\midrule
E1 (commune)   & - & $\checkmark$ & - & $\checkmark$ & - & $\checkmark$ & - & - & $\checkmark$ & - & $\checkmark$ & $\checkmark$ & - \\
E2 discipl. IE & $\checkmark$ & - & $\checkmark$ & - & $\checkmark$ & - & $\checkmark$ & $\checkmark$ & - & - & $\checkmark$ & $\checkmark$ & $\checkmark$ \\
\bottomrule
\end{tabular}
\end{center}
\smallskip
\footnotesize\textbf{Sous-thèmes fréquents :}
MLI vs pleine onde - THD (6/13) - identification convertisseurs chaîne (7/13) - onduleur triphasé + machine (5/13) - calcul de $f_s$ (3/13) - îlotage vs réseau (3/13)

\medskip
\textbf{Ordres de grandeur (repères rapides).}
\begin{itemize}[leftmargin=1.3em]
  \item Réseau BT FR : tension phase-neutre typique \SI{230}{\volt}, tension phase-phase typique \SI{400}{\volt}, fréquence \SI{50}{\hertz}.
  \item Courants d'étude fréquents en sujet : de quelques \si{\ampere} à quelques dizaines de \si{\ampere}.
  \item En conversion d'énergie, un rendement attendu reste généralement entre 0{,}8 et 0{,}98 selon la technologie.
\end{itemize}

\medskip
\textbf{Synthèse express (à refaire sans support).}
\begin{enumerate}[leftmargin=1.4em]
  \item Poser les hypothèses du modèle et les conventions de signe.
  \item Écrire 2 à 4 relations clés avec unités.
  \item Réaliser un calcul numérique complet et vérifier l'ordre de grandeur.
  \item Justifier le résultat par un argument technique de niveau CAPET.
\end{enumerate}

\end{tcolorbox}

\CAPETENDFILE
